\section{The Vocation}

\begin{quotex}
But the sensual, or natural, man perceives not these things that are of the Spirit of God; for it is foolishness to him, and he cannot understand them, because they are spiritually discerned. But the spiritual man judges all things; and he himself is judged by no man.
\flright{\textsc{1 Corinthians 2:14-15}}
\end{quotex}


Those verses show that distinctions between different types of men has scriptural roots. The lower cannot judge the higher, obviously because they can't understand it. Yet that does not deter them. The more intelligent of the natural men (i.e., those dominated by the intellectual center) prefer debates and discussions; they are always on the lookout for the next book and seldom can be persuaded to actually engage in self-development or self-purification.

They share opinions on social media, not so much to exchange ideas, but rather to demonstrate ``gang colors''. Those holding similar opinions, read the same books, or valorize the same public figures, form a common gang. \textbf{Miguel de Molinos}, in \emph{The Spiritual Guide}, has a deeper and more correct understanding:

\begin{quotationx}
Mystical knowledge proceeds not from Wit, but from Experience; it is not invented, but proved; not read, but received; and is therefore most secure and efficacious, of great help and plentiful in fruit; it enters not into the soul by ears, nor by the continual reading of books, but by the free Infusion of the Holy Ghost, whose Grace with most delightful intimacy, is communicated to the simple and lowly.

Besides, it is certain, that he who hath not the experience of this sweetness, cannot pass a Judgement upon these Mysterious Secrets; nay, rather he'll be Scandalized (as many are) when he hears of the Wonders which the Divine Love is wont to work in Souls, because he finds no such Rarities in his own. This science is not theoretical, but practical, wherein experience surpasses the most refined and ingenious speculation. It is to be taken notice of, that the doctrine of this book instructs not all sorts of persons, but those only who have the senses and passions well mortified, who have already advanced and made progress in prayer, and are called by God to the inward way, who encourages and guides them, freeing them from the obstacles which hinder the course to perfect Contemplation.
\end{quotationx}


Social media provides merely a shallow form of friendship, one of convenience or pleasure. Spiritual friendship\footnote{\url{https://www.meditationsonthetarot.com/spiritual-friendship}} is quite different. A true friend is like another ``self'', since he has the same character, shares in common activities, and works toward the same goals. As \textbf{Valentin Tomberg} points out in the Meditations,

\begin{quotex}
Spiritual exercises in common form the common link that unites Hermetists. It is not knowledge in common which unites them, but rather the spiritual exercises and the experience which goes hand in hand with them. That which one knows is the result of personal experience and orientation, whilst \emph{depth}, the level to which one attains --- disregarding the aspect and extent of knowledge that one has gained --- is what is in common.
\end{quotex}


Hermetists, therefore, will have different intellectual interests, vocations, and destinies. But spiritual exercises link them together. Hence, in private groups and esoteric schools, exercises and practices take precedence over erudite displays of knowledge. It helps to practice together and share experiences of the work. This fulfills \textbf{E. F. Schumacher}'s four fields of knowledge:

\begin{enumerate}


\item Knowledge of self

\item Knowledge of the inner life of other people

\item Knowledge of how others experience you, i.e., external considering

\item Knowledge of the world, through physics and metaphysics
\end{enumerate}


Through sharing experiences with spiritual friends, one can develop proficiency in the four fields in a supportive atmosphere; tasks are assigned, just as a scientist designs controlled experiments where theories can be tested. In \emph{Initiation and Spiritual Realization}, \textbf{Rene Guenon} points out the value of collective work:

\begin{quotex}
When the sages converse about the divine mysteries, the Shekinah is present among them ... 
\end{quotex}


Guenon mentions practices like self-purification, painful efforts, sacrifice, detachment and so on. These practices are never mentioned on social media, which instead focus on metaphysics, history of thought, boasting about books, etc., or worse, apotheosis, as though praise causes self-development. Unfortunately, many hold the mistaken belief that Guenon is introducing a new doctrine of founded a new ``school'', whereas he is just restating traditional doctrines. To understand Guenon, first understand a particular tradition as well as its practices. That is why Guenon emphases the necessity of participating in an exoteric religion.

\textbf{Aristotle} said that each being must accomplish his proper activity, which requires two conditions:

\begin{enumerate}
\item The activity is in conformity to his own nature

\item The activity actualizes the being's possibilities
\end{enumerate}


Thus, for any work to be effective, it must correspond to the person's vocation. This vocation is not for everyone.

Notes from group meeting of 26 Oct 2020


\flrightit{Posted on 2020-10-30 by Cologero}

